\chapter{Introduction and Planning}
\label{chapter:introduccion_and_planning}

\section{Introduction}

This Final Master’s Project focuses on designing and implementing a distributed C++ program for simulating flood events. Floods are 
sudden and violent events that have serious economic and human consequences. Specifically, this project was launched in response to 
the DANA (Depresión Aislada en Niveles Altos) that struck Valencia, Spain, on October 29, 2024. This was not just any cold snap, but the most catastrophic one in the 
country’s recent history. For this reason, it has been used as a case study to test modeling capabilities.

The proposed system uses a Cellular Automata engine to model water flow. The terrain is divided into a grid of cells with topographic 
data from Geographic Information Systems. The simulator applies rules to calculate changes in water depth and velocity over time. 
This approach allows for processing of large-scale Digital Elevation Models while maintaining precision.

A key feature of this project is its distributed architecture. The system operates in an environment using the MQTT protocol for 
communication between the simulation engine and external visualization modules. 

This structure is a step towards creating a Digital Twin for disaster management. By integrating real-time data feeds with the 
Cellular Automata model the system aims to provide a synchronized representation of the flood. The project aims to bridge the gap 
between flood modeling and practical emergency response.

By simulating the Valencia flood the work shows how distributed computing and cellular modeling can be used to identify high-risk 
zones and predict flood evolution in time. This provides a tool, for both analyzing past floods and mitigating future crises.


\section{Motivation and Justification}

The increasing frequency of extreme weather events poses a severe threat to urban infrastructure and emergency services. The catastrophic 
flood in Valencia on October 29, 2024, exposed the limitations of traditional predictive frameworks, highlighting the critical need for 
rapid, real-time simulation tools. To address this challenge, this research proposes an efficient, interconnected framework built upon 
three pillars:

\begin{itemize}
    \item \textbf{Technical Efficiency (Cellular Automata)}: Traditional hydrodynamic models are computationally expensive and ill-suited 
    for real-time forecasting. This project leverages Cellular Automata (CA) to provide a streamlined, highly parallelizable alternative 
    for efficient fluid routing.

    \item \textbf{Distributed Innovation (MQTT Integration)}: To bridge the gap between simulation and emergency response, the framework 
    integrates the MQTT protocol into a distributed architecture. This enables real-time 3D geospatial visualization, scalable field-sensor 
    integration, and the creation of a dynamic Digital Twin for predictive risk management.

    \item \textbf{Socio-Economic Impact}: Beyond algorithmic innovation, the system provides immediate utility for civil protection by 
    enabling post-event forensic analysis to identify structural vulnerabilities, as well as real-time decision support for optimized 
    evacuations and resource deployment.
\end{itemize}


\section{Objectives}

\subsection{General Objective}

The main goal is to create a system that can simulate phenomena such as floods. It will use Cellular Automata and the C++ programming 
language to make it work well. The project wants to provide a way to analyze disasters and monitor them in time. This system will be 
like a Digital Twin.

\subsection{Specific Objectives}

To achieve this goal, the following specific objectives have been established:

\begin{itemize}
    \item \textbf{Literature Review}: Analyze existing scientific literature on Cellular Automata (CA) and C++ frameworks applied to hydrodynamic and flood modeling.
    \item \textbf{CA Engine Design}: Define mathematical transition rules and neighborhood topologies (e.g., Moore or Von Neumann) to simulate water depth dynamics.
    \item \textbf{GIS Data Integration}: Develop a parsing module for Digital Elevation Models (DEM) to accurately initialize the environment with real-world topography.
    \item \textbf{Distributed Communication}: Implement an MQTT-based protocol layer to enable real-time data exchange with external visualizers within a distributed network.
    \item \textbf{Model Validation}: Validate the core water propagation logic by bench-testing simulation outputs against empirical data from the 2024 Valencia DANA event.
    \item \textbf{Visualization and Post-Processing}: Create post-processing utilities to interpret, map, and analyze the flood's spatial impact on the terrain.
\end{itemize}


\section{Project Planning and Management}
\label{sec:project_planning_and_management}

The successful execution of the Final Master's Project (FMP), which focuses on developing a scientific C++ program for flood simulation, 
requires robust planning and diligent management. This course is weighted at 18 ECTS, equating to an estimated dedicated effort of 
approximately 450 hours. The timeline presented below serves as a guiding roadmap, designed to track progress against key project 
milestones and deliverables.

\subsection{Work Methodology}
\label{subsec:methodology}

The project methodology is structured around the four mandatory phases defined for the FMP, integrated with a phased software 
development approach. This hybrid methodology ensures that the core theoretical and design elements (Phases 1 \& 2) are completed 
before critical implementation and validation (Phase 3) and final documentation (Phase 4). The iterative nature of the plan allows 
for necessary adjustments and includes dedicated time for thorough testing and writing throughout the academic year.

\subsection{Project Schedule: Gantt Chart}
\label{subsec:gantt}

The following Gantt chart (Figure \ref{fig:ganttchart}) illustrates the proposed 24-week timeline for the project, distributing the 
estimated 450 working hours across the defined phases. Time buffers have been implicitly integrated into the task durations to account 
for unforeseen issues, necessary code refactoring, and academic holidays.

% ----------------------------------------------------------------
% START OF GANTT CHART CODE
% ----------------------------------------------------------------

\begin{figure}[h!]
    \resizebox{\linewidth}{!}{
        \centering
        \begin{ganttchart}[
            vgrid,
            bar height=0.6,
            group height=0.3,
            title label font=\bfseries\normalsize
        ]{1}{24}
            \gantttitle{Project Timeline (24 Weeks)}{24} \\
            \gantttitlelist{1,...,24}{1} \\
            % PHASE 1
            \ganttgroup[group label node/.append style={fill=blue!20}]{Phase 1: Planning (4 Weeks)}{1}{4} \\
            \ganttbar[bar/.append style={fill=blue!50}]{P1.1 State-of-the-Art Review}{1}{3} \\
            \ganttbar[bar/.append style={fill=blue!50}]{P1.2 Objectives and Scope Definition}{3}{4} \\
            \ganttbar[bar/.append style={fill=blue!50}]{P1.3 Development Environment Setup}{1}{1} \\ % Etiqueta acortada
            % PHASE 2
            \ganttgroup[group label node/.append style={fill=green!20}]{Phase 2: Problem \& Proposal (6 Weeks)}{5}{10} \\
            \ganttbar[bar/.append style={fill=green!50}]{P2.1 Cellular Automata (CA) Design}{5}{6} \\
            \ganttbar[bar/.append style={fill=green!50}]{P2.2 GIS Data Reading Implementation}{5}{7} \\ % Etiqueta acortada
            \ganttbar[bar/.append style={fill=green!50}]{P2.3 Water Propagation Rules (CA)}{7}{10} \\
            % PHASE 3
            \ganttgroup[group label node/.append style={fill=yellow!20}]{Phase 3: Partial Deliverable (6 Weeks)}{11}{16} \\
            \ganttbar[bar/.append style={fill=yellow!70}]{P3.1 Unit Testing of the CA Engine}{11}{12} \\
            \ganttbar[bar/.append style={fill=yellow!70}]{P3.2 Integration and Data Validation}{13}{15} \\
            \ganttbar[bar/.append style={fill=yellow!70}]{P3.3 Progress Document Generation. (Part 1)}{14}{16} \\ % Etiqueta acortada
            % PHASE 4
            \ganttgroup[group label node/.append style={fill=red!20}]{Phase 4: Final Deliverable (6 Weeks)}{17}{22} \\
            \ganttbar[bar/.append style={fill=red!50}]{P4.1 Results Analysis and Visualization}{17}{19} \\
            \ganttbar[bar/.append style={fill=red!50}]{P4.2 Final Thesis Writing and Revision}{18}{21} \\
            \ganttbar[bar/.append style={fill=red!50}]{P4.3 Oral Presentation Preparation}{20}{22} \\
            % PHASE 5
            \ganttgroup[group label node/.append style={fill=gray!20}]{Phase 5: Presentation and Defense (2 Weeks)}{23}{24} \\
            \ganttbar[bar/.append style={fill=gray!50}]{P5.1 Final Supervisor Review}{23}{23} \\
            \ganttbar[bar/.append style={fill=gray!50}]{P5.2 Virtual Defense}{24}{24}
        \end{ganttchart}
    }
    \caption{Project Timeline (24 Weeks)}
    \label{fig:ganttchart}
\end{figure}

% ----------------------------------------------------------------
% END OF GANTT CHART CODE
% ----------------------------------------------------------------

\subsection{Details of Activities}
\label{subsec:activity_details}

The following table (Table \ref{tab:fmp_activities}) provides a detailed breakdown of the tasks defined in the timeline, specifying the scope of work for each activity.

% ----------------------------------------------------------------
% START OF LONGTABLE CODE
% ----------------------------------------------------------------

\begin{table}[htbp]
\centering
\small % Reduce un poco el tamaño de la fuente para asegurar el ajuste
\caption{Detailed Breakdown of Final Master's Project Activities.}
\label{tab:fmp_activities}
\begin{tabularx}{\textwidth}{l l X}
    \toprule
    \textbf{ID} & \textbf{Activity} & \textbf{Description} \\
    \midrule

    % PHASE 1
    \multicolumn{3}{l}{\textbf{PHASE 1: PLANNING}} \\
    \midrule
    P1.1 & State-of-the-Art Review & Review literature on flood simulation, focusing on Cellular Automata (CA) and C++ flow models. \\
    P1.2 & Objectives \& Scope & Define flood scenarios (2D/3D), terrain types, success criteria, and required GIS input formats. \\
    P1.3 & Environment Setup & Configure C++ compiler, GIS libraries (GDAL), Git, and the initial LaTeX document structure. \\
    \midrule

    % PHASE 2
    \multicolumn{3}{l}{\textbf{PHASE 2: PROBLEM DESCRIPTION AND PROPOSAL}} \\
    \midrule
    P2.1 & CA Design & Define transition rules, neighborhood structure, and C++ structures for terrain and water depth. \\
    P2.2 & GIS Data Reading & Implement C++ module to read DEM data and initialize CA states based on topography. \\
    P2.3 & Water Propagation Rules & Code the simulator core with water transfer equations governing the flood simulation iterations. \\
    \midrule

    % PHASE 3
    \multicolumn{3}{l}{\textbf{PHASE 3: PARTIAL DELIVERABLE}} \\
    \midrule
    P3.1 & Unit Testing & Test critical functions (GIS reading, CA rules) to fix precision and performance errors. \\
    P3.2 & Integration \& Validation & Run simulator with representative GIS data and validate results against known flood data. \\
    P3.3 & Progress Document (Part 1) & Draft LaTeX chapters: introduction, literature review, methodology, and CA design. \\
    \midrule

    % PHASE 4
    \multicolumn{3}{l}{\textbf{PHASE 4: FINAL DELIVERABLE}} \\
    \midrule
    P4.1 & Results \& Visualization & Process output data (depth maps, time) and generate graphs to interpret flood impact. \\
    P4.2 & Thesis Writing \& Revision & Complete Results/Conclusions chapters in LaTeX and perform full grammatical review. \\
    P4.3 & Oral Presentation Prep & Design slides and rehearse speech, emphasizing CA model, C++ code, and GIS results. \\
    \midrule

    % PHASE 5
    \multicolumn{3}{l}{\textbf{PHASE 5: PRESENTATION AND VIRTUAL DEFENSE}} \\
    \midrule
    P5.1 & Supervisor Review & Final session with supervisor to apply last-minute corrections to document, code, and slides. \\
    P5.2 & Virtual Defense & Present FMP to the evaluation committee, defending technical decisions and conclusions. \\
    \bottomrule
\end{tabularx}
\end{table}

% ----------------------------------------------------------------
% END OF LONGTABLE CODE
% ----------------------------------------------------------------


